% Provenance: paper_P4_length_bias + sections/length_bias.tex (incl. scale
% extension), E-P4 corrected panel artifacts (p4uncap_*.json), cell_runner loss
% arms, min_report loss-form item.
\chapter{Loss Form and Length Bias}
\label{ch:lossform}

\section{The question}
Dr.GRPO's motivating thesis is that GRPO's per-completion normalisations bias
updates toward longer outputs --- the ``verbosity trap'': length inflates,
then reward collapses. This chapter tests for the signature at our scale with
per-step traces of mean reward and mean completion length.

\section{Short-horizon panels at 0.5B--1.5B}
On Qwen2.5-0.5B/arithmetic (40 steps, 5 seeds per loss) and
Qwen2.5-1.5B-Instruct/GSM8K-CoT (30 steps, 3 seeds per loss), neither GRPO nor
Dr.GRPO exhibits the trap: per-run $\rho(\text{step},\text{len})$ is
consistently \emph{negative} (models get shorter as training proceeds). These
runs are cap-bounded (200-token generation cap) and short-horizon; the
original Dr.GRPO evidence sits at hundreds of steps, an order of magnitude
beyond ours --- our data is consistent with the early phase of that curve.

\section{The uncapped 8B panel}
A six-arm panel extends the scale axis: Qwen3-8B, GSM8K, 30 steps, $G{=}8$,
batch 4, \textbf{three seeds per loss}, completion cap raised $5\times$ to
1{,}024 tokens. Results:
\begin{center}
\begin{tabular}{@{}lccc@{}}
\toprule
arm & len first-5 $\to$ last-10 & late \ZVF{} & late reward\\
\midrule
GRPO s42/s123/s456 & $1004{\to}905$, $981{\to}944$, $996{\to}900$ &
$0.47$, $0.47$, $0.55$ & $0.75$, $0.68$, $0.76$\\
Dr.GRPO s42/s123/s456 & $999{\to}931$, $972{\to}902$, $1000{\to}878$ &
$0.45$, $0.72$, $0.70$ & $0.74$, $0.78$, $0.78$\\
\bottomrule
\end{tabular}
\end{center}
\textbf{No length inflation in either loss} --- all six arms shrink 6--12\%
--- and \textbf{no late-\ZVF{} separation} between the losses. Step-0 lengths
sit near the raised cap ($\approx98\%$), so upward headroom remains limited,
but trajectories move \emph{away} from the cap, so the compression finding is
not a censoring artifact. The 30-step horizon caveat stands.

\section{Provenance: this panel was run twice}
The first version of this panel was \emph{invalid}: the runner's
\texttt{--loss drgrpo} flag was documented (its commit message claimed the
mapping) but never wired to the advantage normalisation, so both ``arms''
silently trained vanilla GRPO --- and no reward, length, or \ZVF{} trace
revealed it. The defect was caught only by reading the runner. The invalid
artifacts are preserved under explicit \texttt{.invalid\_actually\_grpo.json}
names; the corrected rerun (same day) produced the table above, and notably
\emph{reversed} the invalid panel's secondary claim (the invalid data had
suggested Dr.GRPO keeps more groups live; the corrected data shows
equal-or-higher late \ZVF{} for Dr.GRPO). Chapter~9 turns this incident into
protocol.

\section{Conclusion of the chapter}
At the scales and horizons this programme can afford, the loss-form choice
between GRPO and Dr.GRPO has no observable footprint on length or \ZVF{}.
This is evidence about \emph{reporting} rather than about superiority: a
comparison between these losses at this scale measures noise unless the stack
is controlled far more tightly than the labels suggest.
